\section{Formalization in Lean~4}
\label{sec:formalization}

\subsection{Project structure and statistics}
\label{sec:project-structure}

The formalization is organized into 41~Lean source files, grouped by track
in \cref{tab:files}. Each file is self-contained with explicit imports,
enabling parallel development and independent verification.

\begin{table}[ht]
\centering
\footnotesize
\setlength{\tabcolsep}{4pt}
\begin{tabular}{@{}lrl@{}}
\toprule
\textbf{File} & \textbf{Lines} & \textbf{Content} \\
\midrule
\multicolumn{3}{@{}l}{\textit{Lie--Trotter core (Tracks 1--4, \cref{sec:proof})}} \\
\leansrc{Telescoping.lean}       & 106  & Algebraic telescoping + norm bound \\
\leansrc{ExpBounds.lean}         & 212  & Exponential series remainder estimates \\
\leansrc{StepError.lean}         & 674  & Quadratic step error (algebraic factorization) \\
\leansrc{ExpDivPow.lean}         & 54   & $(e^{a/n})^n = e^a$ \\
\leansrc{Assembly.lean}          & 146  & $O(1/n)$ convergence + main theorem \\
\midrule
\multicolumn{3}{@{}l}{\textit{Strang, multi-operator, commutator scaling (Track~5)}} \\
\leansrc{StrangSplitting.lean}   & 492  & Strang splitting, $O(1/n^2)$ \\
\leansrc{MultiOperator.lean}     & 288  & Multi-operator first-order \\
\leansrc{MultiStrang.lean}       & 428  & Multi-operator Strang \\
\leansrc{CommutatorScaling.lean} & 379  & Duhamel commutator-scaling \\
\leansrc{MultiCommutatorScaling.lean}       & 128  & Multi-operator comm.\ scaling \\
\leansrc{StrangCommutatorScaling.lean}      & 733  & Strang double-commutator scaling \\
\leansrc{MultiStrangCommutatorScaling.lean} & 174  & Multi-op.\ Strang comm.\ scaling \\
\leansrc{HigherCommutator.lean}            & 243  & Triple-FTC: extracts $[B,[B,[B,A]]]$ \\
\leansrc{StrangCommutatorScalingTight.lean} & 603 & Tighter Strang via norm-of-difference $D$; $\norm{D}$ domination \\
\leansrc{StrangTotalErrorCommScaling.lean}  & 184 & Commutator-scaled Strang \emph{total} error \\
\midrule
\multicolumn{3}{@{}l}{\textit{Suzuki $S_4$ track (\cref{sec:suzuki4})}} \\
\leansrc{Suzuki4.lean}                    & 466  & $S_4$ definition; generic $O(1/n^2)$ rate \\
\leansrc{Suzuki4FullDuhamel.lean}         & 445  & $S_4$ $O(t^3)$ via 5-$S_2$ telescoping \\
\leansrc{Suzuki4CommutatorScaling.lean}   & 90   & \texttt{suzuki4Exp} definition \\
\leansrc{Suzuki4Deriv.lean}               & 96   & Derivative scaffolding \\
\leansrc{Suzuki4HasDerivAt.lean}          & 136  & Module~1: \texttt{HasDerivAt} for 12-factor $w_4$ \\
\leansrc{Suzuki4Module2.lean}             & 167  & Module~2: FTC-2 bridge \\
\leansrc{Suzuki4Module3.lean}             & 184  & Module~3: FTC-2 reduction \\
\leansrc{Suzuki4Module4.lean}             & 175  & Module~4a: continuity of $w_4'$ \\
\leansrc{Suzuki4DerivExplicit.lean}       & 979  & Module~4b: explicit derivative; Phase 1--3 \\
\leansrc{Suzuki4StrangBlocks.lean}        & 120  & $S_4$ = 5 Strang blocks + cubic cancellation \\
\leansrc{Suzuki4MultinomialExpand.lean}   & 787  & Multinomial formulas; h2 unconditional, h3 under Suzuki cubic \\
\leansrc{Suzuki4Phase5.lean}              & 782  & Taylor reduction + Leibniz bridges + CAPSTONE \\
\leansrc{Suzuki4OrderFive.lean}           & 440  & $O(t^5)$ abstract-form target \\
\leansrc{Suzuki4ChildsForm.lean}          & 216  & Childs Prop pf4\_bound\_2term framework \\
\leansrc{Suzuki4BchBound.lean}            & 495  & SLICE~1: single-step $O(\abs{\tau}^5)$ bound \\
\leansrc{TaylorMatch.lean}                & 331  & SLICE~2: generic Taylor-match-from-norm \\
\leansrc{Suzuki4ViaBCH.lean}              & 1\,436 & SLICE~3 wiring + L1/L2/L3/L4 BCH-derived bounds \\
\leansrc{Suzuki4Convergence.lean}         & 398  & Total error $O(1/n^4)$ + convergence \\
\leansrc{Suzuki4TightConvergence.lean}    & 169  & Commutator-scaled $O(1/n^4)$ \\
\leansrc{Suzuki4UnitaryTotalError.lean}   & 274  & Skew-adjoint total error, no growth factor \\
\leansrc{Suzuki4Commute.lean}             & 228  & Commuting case: $S_4$ exact, error $\equiv 0$ \\
\leansrc{Suzuki4GapClosers.lean}          & 205  & General-$t$ step; imaginary-time corollaries \\
\leansrc{TrotterStepCount.lean}           & 392  & $\varepsilon$-form step counts; explicit per-$n$ bounds \\
\leansrc{MatrixCorollaries.lean}          & 157  & Matrix specializations (spectral norm) \\
\leansrc{PrefactorStrict.lean}            & 111  & Strict $\gamma_i < \alpha_i$; gap $\gamma_i \le \alpha_i/8$ \\
\midrule
\multicolumn{3}{@{}l}{\textit{Root}} \\
\href{\repourl/LieTrotter.lean}{\texttt{LieTrotter.lean}} & 73 & Import aggregator \\
\midrule
\textbf{Total} & \textbf{14\,196} & \textbf{0 \texttt{sorry} axioms,
0 own theorem-level BCH-interface axioms} \\
\bottomrule
\end{tabular}
\caption{Source files in the formalization, grouped by track.
\texttt{Suzuki4.lean} carries the generic $O(1/n^2)$ rate, which holds for
every $n \ge 1$ and needs no cubic condition;
\texttt{Suzuki4Convergence.lean} upgrades the \emph{same} object to $O(1/n^4)$
under the Suzuki cubic condition (\cref{sec:s4-convergence}), and
\texttt{Suzuki4TightConvergence.lean} does so with a commutator-scaled leading
coefficient.}
\label{tab:files}
\end{table}

The development builds against Lean~4.29.0-rc8 and current \Mathlib.
The companion
\href{https://github.com/Jue-Xu/Lean-BCH}{Lean-BCH} project is imported
as a git dependency, supplying the symmetric three-factor BCH identities,
the two-factor palindromic quintic remainder, and the five-factor
palindromic quintic identification used in the Suzuki~$S_4$ track.
Lean-BCH is sorry-free as of revision~\texttt{d455ff0} (2026-05-19),
which closed the previous B1.c (\texttt{symmetric\_bch\_quintic\_sub\_poly\_axiom})
dependency; its single-step $O(\abs{\tau}^5)$ payoff lemma anchors
our SLICE-1 bound in \leansrc{Suzuki4BchBound.lean}. Two private
septic axioms remain in Lean-BCH and gate only the optional Level~4
local small-$\tau$ ($\tau^7$) refinement (see \cref{sec:caveats}).
Build time for the full project is under 2~minutes on a standard
laptop after \Mathlib oleans are cached (\texttt{lake exe cache get}).

\subsection{Type-level setup}
\label{sec:type-setup}

Our theorems are stated for a general complete normed algebra $\A$ over a
field $\mathbb{K} \in \{\mathbb{R}, \mathbb{C}\}$:

\begin{leancode}
variable {𝕂 : Type*} [RCLike 𝕂]
variable {𝔸 : Type*} [NormedRing 𝔸] [NormedAlgebra 𝕂 𝔸]
  [NormOneClass 𝔸] [CompleteSpace 𝔸]
\end{leancode}

The typeclass \texttt{NormOneClass} requires $\norm{1} = 1$, which is needed
for \texttt{norm\_pow\_le} in current \Mathlib. The main theorem is
(\leanref{Assembly.lean}{121}{lie\_trotter}):

\begin{leancode}
theorem lie_trotter (A B : 𝔸) :
    Filter.Tendsto
      (fun n : ℕ => (exp ((n : 𝕂)⁻¹ • A) *
                      exp ((n : 𝕂)⁻¹ • B)) ^ n)
      atTop (nhds (exp (A + B)))
\end{leancode}

Note the scalar action $\bullet$: the step is $n^{-1}$ \emph{acting on} $A$
in the $\mathbb{K}$-module $\A$, not a ring product with a coerced $n^{-1} \in
\A$. The two coincide in a unital algebra but are distinct terms in Lean, and
the derivative and integral lemmas we rely on are all stated for $\bullet$.

\paragraph{The \texttt{include} pattern.}
Since \texttt{NormedSpace.exp} no longer takes a field parameter in recent
\Mathlib, the variable $\mathbb{K}$ does not appear in the type of theorems
involving \texttt{exp}---in \texttt{norm\_exp\_le} below, both sides mention
only $\A$ and $\mathbb{R}$. Lean~4 drops unused type-level variables, so we
must write \texttt{include $\mathbb{K}$ in} before each lemma whose proof (but
not statement) requires $\mathbb{K}$
(e.g., \leanref{ExpBounds.lean}{72}{norm\_exp\_le}, whose proof rewrites with
\texttt{exp\_tsum\_form ($\mathbb{K}$ := $\mathbb{K}$)}):

\begin{leancode}
include 𝕂 in
lemma norm_exp_le (a : 𝔸) : ‖exp a‖ ≤ Real.exp ‖a‖ := by ...
\end{leancode}

\subsection{Telescoping in Lean}
\label{sec:lean-telescoping}

The telescoping identity~\eqref{eq:telescoping} is proved by induction on $n$
using \texttt{Finset.sum\_range\_succ} to peel off the last term
(\leanref{Telescoping.lean}{32}{telescoping\_direct}):

\begin{leancode}
theorem telescoping_direct (X Y : R) (n : ℕ) :
    (∑ k ∈ Finset.range n,
      X ^ k * (X - Y) * Y ^ (n - 1 - k)) =
      X ^ n - Y ^ n := by
  induction n with
  | zero => simp
  | succ m ih =>
    rw [Finset.sum_range_succ]
    -- Factor Y from the inner sum to get the
    -- inductive hypothesis
    ...
    noncomm_ring [pow_succ]
\end{leancode}

Here \texttt{R} is any \texttt{Ring}: the telescoping identity is purely
algebraic, so this file assumes no norm and no completeness.

The final step uses \texttt{noncomm\_ring}, which handles polynomial identities
in non-commutative rings. The standard \texttt{ring} tactic silently fails on
non-commutative goals---it produces an unsolved subgoal rather than an error,
which cost us significant debugging time early in the project.

\subsection{Exponential bounds: filling Mathlib gaps}
\label{sec:lean-exp-bounds}

The bound $\norm{e^a} \le e^{\norm{a}}$ was not available in \Mathlib for
general Banach algebras. Our proof works from the tsum representation
(\leanref{ExpBounds.lean}{72}{norm\_exp\_le}):

\begin{leancode}
include 𝕂 in
lemma norm_exp_le (a : 𝔸) : ‖exp a‖ ≤ Real.exp ‖a‖ := by
  rw [exp_tsum_form (𝕂 := 𝕂), real_exp_eq_tsum]
  calc ‖∑' n, (((Nat.factorial n : 𝕂)⁻¹) • a ^ n : 𝔸)‖
      ≤ ∑' n, ‖...‖         -- norm_tsum_le_tsum_norm
    _ ≤ ∑' n, ‖a‖ ^ n / n ! -- norm_exp_term_le
    _ = Real.exp ‖a‖        -- Real.hasSum_exp
\end{leancode}

The proof chains three standard ingredients: pulling the norm inside the tsum,
bounding each term $\norm{a^n/n!} \le \norm{a}^n/n!$ via submultiplicativity,
and recognizing the resulting series as $e^{\norm{a}}$.

\subsection{Step error: algebraic factorization in Lean}
\label{sec:lean-step-error}

The factorization~\eqref{eq:factorization} is verified by \texttt{noncomm\_ring}.
The cross-term bound~\eqref{eq:cross-term} is proved by induction using the
non-commutative recurrence
(\leanref{StepError.lean}{43}{norm\_pow\_add\_sub\_pow\_sub\_pow}):

\begin{leancode}
private lemma norm_pow_add_sub_pow_sub_pow (a b : 𝔸) :
    ∀ m : ℕ, 1 ≤ m →
      ‖(a + b) ^ m - a ^ m - b ^ m‖ ≤
        (‖a‖ + ‖b‖) ^ m - ‖a‖ ^ m - ‖b‖ ^ m := by
  intro m hm
  induction m with
  | zero => omega
  | succ n ih => ...
\end{leancode}

At 674~lines, \leansrc{StepError.lean} is the longest file in the basic
Lie--Trotter development (excluding commutator-scaling extensions).
The inductive cross-term bound accounts for roughly half of this length,
because each step requires explicit \texttt{norm\_mul\_le} applications and
\texttt{gcongr} for the inductive comparison.

\subsection{Commutator scaling: calculus in Lean}
\label{sec:lean-commutator}

The commutator-scaling proofs are the most technically demanding part of the
formalization, as they require doing analysis---derivatives, interval integrals,
FTC-2---in a non-commutative Banach algebra.

\paragraph{Derivative computations.}
\Mathlib's \texttt{hasDerivAt\_exp\_smul\_const'} supplies the
derivative $d/d\tau[e^{\tau B}] = B\,e^{\tau B}$, and the
product rule \texttt{HasDerivAt.mul} extends to normed algebras.
One subtlety: $(-\tau) \bullet B$ and $\tau \bullet (-B)$ are
definitionally distinct but equal; we must normalize before
applying derivative lemmas
(\leanref{CommutatorScaling.lean}{75}{hasDerivAt\_exp\_conj}):

\begin{leancode}
lemma hasDerivAt_exp_conj (A B : 𝔸) (s : ℝ) :
    HasDerivAt
      (fun u => exp (u • B) * A * exp ((-u) • B))
      (exp (s • B) * (B * A - A * B) *
       exp ((-s) • B)) s := by
  simp_rw [show ∀ u : ℝ, (-u) • B = u • (-B)
    from fun u => by rw [neg_smul, smul_neg]]
  ...
\end{leancode}

\paragraph{FTC-2 for integral representations.}
The conjugation identity~\eqref{eq:conj-ftc} and the integral error
representation~\eqref{eq:integral-error} are proved by
\texttt{integral\_eq\_sub\_of\_hasDerivAt}
(\leanref{CommutatorScaling.lean}{157}{exp\_conj\_sub\_eq\_integral}):

\begin{leancode}
theorem exp_conj_sub_eq_integral (A B : 𝔸) (τ : ℝ) :
    exp (τ • B) * A * exp ((-τ) • B) - A =
      ∫ s in (0 : ℝ)..τ,
        exp (s • B) * (B * A - A * B) *
        exp ((-s) • B) := by
  have hderiv : ∀ u ∈ Set.uIcc 0 τ,
      HasDerivAt ... := fun u _ =>
    hasDerivAt_exp_conj A B u
  have hint : IntervalIntegrable ... :=
    (continuous_exp_conj_deriv A B).intervalIntegrable 0 τ
  exact (integral_eq_sub_of_hasDerivAt hderiv hint).symm
\end{leancode}

\paragraph{Pulling constants through integrals.}
In a \texttt{NormedRing}, left-multiplication by a fixed element $c$ is
\emph{not} a scalar multiplication---it is ring multiplication. To derive
$c \cdot \int f = \int c \cdot f$, we use
\texttt{ContinuousLinearMap.intervalIntegral\_comp\_comm} with
$L = \texttt{ContinuousLinearMap.mul}\ \mathbb{R}\ \A\ c$, which views
left-multiplication as a continuous linear map.

\paragraph{Working over $\mathbb{R}$.}
The commutator-scaling files work over $\texttt{NormedAlgebra}\ \mathbb{R}\ \A$
directly, rather than a general $\mathbb{K}$. The reason is that the
\texttt{HasDerivAt}/\texttt{intervalIntegral} machinery requires
$\texttt{SMul}\ \mathbb{R}\ \A$, which is \emph{not} automatically synthesized
from $\texttt{[RCLike K] [NormedAlgebra K A]}$. For applications with
$\mathbb{K} = \mathbb{C}$, one can use
$\texttt{NormedAlgebra.restrictScalars}\ \mathbb{R}\ \mathbb{C}\ \A$.

\subsection{Multi-operator and Suzuki extensions}
\label{sec:lean-extensions}

The multi-operator formulas are proved by induction on the list of operators.
The palindromic product is defined recursively
(\leanref{MultiStrangCommutatorScaling.lean}{42}{palindromicProd}):

\begin{leancode}
def palindromicProd : List 𝔸 → 𝔸
  | [] => 1
  | [a] => exp a
  | a :: rest =>
    exp ((1/2 : ℝ) • a) *
    palindromicProd rest *
    exp ((1/2 : ℝ) • a)
\end{leancode}

The step size is not a parameter: it is folded into the operators by the
caller, which passes the list $[t \bullet A_1, \ldots, t \bullet A_m]$. This
keeps the recursion on a single argument, so the induction is on the list
alone. The inductive step peels off the outermost $e^{A_1/2}$ factors and
applies the two-operator Strang bound to the pair
$(A_1, A_2 + \cdots + A_m)$, then recurses on the inner product.
